
\section{问题分析}

\subsection{问题一的分析}

% 第1段：问题要求 + 分析思路（3-5句话，先说"问题一要求我们..."，然后"首先...；接着...；然后...；最后..."）
% 示例：问题一要求我们求解给定条件下XXX的XXX、XXX以及XXX。首先，我们可以通过XXX来确定XXX；接着，通过XXX我们可以由XXX推得XXX，进而推得其XXX；然后，基于上述信息我们可以计算XXX，从而得出XXX，将不同时刻不同XXX求和取平均即可得到XXX；最后，再结合XXX以及XXX的计算公式，通过代入XXX等数据我们就可以求出XXX，进而求出XXX以及XXX。

\subsection{问题二的分析}

% 第1段：问题要求 + 模型类型（2-3句话）
% 示例：问题二要求我们对XXX进行优化，使得XXX在达到XXX的条件下，XXX尽可能大。我们将其理解为一个单目标优化问题。

% 第2段：建模思路（3-4句话，说明目标函数、决策变量、约束条件、求解方法）
% 示例：因此我们以XXX最大为目标函数，以XXX、XXX、XXX以及XXX等参数作为决策变量，再根据题意确立多个约束条件，建立起单目标优化模型。由于要求解的决策变量的数量很多，传统的遍历算法显然是行不通的，因此我们可以采用XXX对该单目标优化模型进行求解。

\subsection{问题三的分析}

% 第1段：与问题二的区别 + 简化思路（3-5句话）
% 示例：相较于问题二，问题三中XXX并不统一，这导致了决策变量的数量进一步增多。为了简化模型，我们可以参考问题二中求得的结论，认为在问题三中XXX仍按照XXX的方式进行XXX，又因为XXX具有XXX性，所以我们可以假设XXX，因此新增的决策变量减少为XXX，其余目标函数和约束条件均与第二问相同。

% 第2段：求解方法（1-2句话）
% 示例：在求解时，我们可以在XXX的基础上对XXX采用XXX的方式进行遍历求解。
